\documentclass[a4paper,10pt,fleqn]{article}
\usepackage[margin=1in]{geometry}
\usepackage{amssymb}
\usepackage{adjustbox}
\usepackage{natbib}
\usepackage[colorlinks=true,linkcolor=blue,citecolor=blue,urlcolor=blue]{hyperref}
\usepackage{fuzz}
\usepackage{zed-maths}
\usepackage{zed-proof}
\newdimen\savedleftskip
\title{Foreign Keys in Relational Schemas}
\author{txt2tex example}
\hypersetup{pdftitle={Foreign Keys in Relational Schemas},pdfauthor={txt2tex example},pdfcreator={txt2tex v1.2.0}}
\begin{document}

\maketitle

\begin{zed}[RecipeId, IngredientId, EmployeeId, CuisineTag, IngredientName, RecipeName, EmployeeName, PositionName]\end{zed}

\begin{axdef}
FK : \power RecipeId \rel \power IngredientId \\
FKEE : \power EmployeeId \rel \power EmployeeId
\end{axdef}

\medskip

\section*{Pattern 1: Single-Attribute Primary Key}

\noindent Each relation has exactly one pk line.  txt2tex underlines the attribute in the schema box and emits a PK statement immediately after.

\bigskip

\begin{schema}{Recipe}
recipeId : RecipeId \\
name : RecipeName \\
cuisine : CuisineTag
\end{schema}
\noindent$\mathrm{PK}(\mathrm{Recipe}) = \{recipeId\}$

\smallskip

\begin{schema}{Ingredient}
ingredientId : IngredientId \\
name : IngredientName
\end{schema}
\noindent$\mathrm{PK}(\mathrm{Ingredient}) = \{ingredientId\}$

\smallskip

\medskip

\section*{Pattern 2: Composite Primary Key}

\noindent Two pk lines in one schema body.  Both attributes are underlined; txt2tex collapses them into one PK set: PK(RecipeIngredient) = $\{recipeId, ingredientId\}$.

\bigskip

\begin{schema}{RecipeIngredient}
recipeId : RecipeId \\
ingredientId : IngredientId
\end{schema}
\noindent$\mathrm{PK}(\mathrm{RecipeIngredient}) = \{recipeId, ingredientId\}$

\smallskip

\medskip

\section*{Pattern 3: Single-Attribute Foreign Key}

\noindent A foreign key is a standalone Z predicate.  The binding-expression form maps one source attribute to one target attribute: (SourceRelation, TargetRelation) $\mapsto$ $\{srcAttr \mapsto tgtAttr\}$ $\in$ FK

\bigskip

\noindent RecipeIngredient.recipeId must match a Recipe.recipeId:

\bigskip

\noindent
$(RecipeIngredient, Recipe) \mapsto \{recipeId \mapsto recipeId\} \in FK$


\noindent RecipeIngredient.ingredientId must match an Ingredient.ingredientId:

\bigskip

\noindent
$(RecipeIngredient, Ingredient) \mapsto \{ingredientId \mapsto ingredientId\} \in FK$


\medskip

\section*{Pattern 4: Composite FK and Self-Referential FK}

\noindent Both FK predicates for RecipeIngredient can be stated together. Each is a separate standalone predicate.

\bigskip

\noindent
$(RecipeIngredient, Recipe) \mapsto \{recipeId \mapsto recipeId\} \in FK$


\noindent
$(RecipeIngredient, Ingredient) \mapsto \{ingredientId \mapsto ingredientId\} \in FK$


\noindent Employee carries a managerId that references another row in the same relation. The self-referential FK predicate uses FKEE, declared above.

\bigskip

\begin{schema}{Employee}
employeeId : EmployeeId \\
name : EmployeeName \\
position : PositionName \\
managerId : EmployeeId
\end{schema}
\noindent$\mathrm{PK}(\mathrm{Employee}) = \{employeeId\}$

\smallskip

\noindent Every managerId in Employee must match a valid employeeId in the same relation:

\bigskip

\noindent
$(Employee, Employee) \mapsto \{managerId \mapsto employeeId\} \in FKEE$


\end{document}